\section{TASK 1: DATASET CHARACTERIZATION} \label{sec:task1}

The labelled data comprise \textbf{251 training sessions} and \textbf{108 test
sessions}. Each session is represented by a raw command string and a list of
MITRE tactics aligned with its space-separated Bash words.

\subsection{Label Distribution} \label{subsec:task1_labels}

\textbf{Distinct tactics and imbalance.} The dataset contains seven tactics:
Discovery, Execution, Persistence, Impact, Defense Evasion, Not Malicious Yet,
and Other. Figure~\ref{fig:task1_tag_distribution} reports their word-level
frequencies in both partitions. The two distributions have a similar shape:
Discovery dominates with \textbf{52.37\%} of training words and
\textbf{53.36\%} of test words, followed by Execution
(\textbf{28.23\%}/\textbf{25.30\%}) and Persistence
(\textbf{9.87\%}/\textbf{11.02\%}). Every remaining class accounts for less
than 4\% of either split.

This imbalance is important for the modelling stage: accuracy alone will be
insufficient, because a classifier can perform well on the dominant tactics
while neglecting minority behaviors such as Impact or Other. Macro-averaged
metrics and per-class F1 scores are therefore necessary.

\begin{figure}[tbp]
    \centering
    \includegraphics[width=0.88\linewidth]{\imagespath Task1/tag_distribution_barplot.png}
    \caption{Distribution of MITRE tactic labels over Bash words in the training
    and test partitions.}
    \label{fig:task1_tag_distribution}
\end{figure}

\subsection{Contextual Meaning of \texttt{echo}} \label{subsec:task1_echo}

\textbf{Observed labels.} In the training set, \texttt{echo} occurs with six
different tactics: Persistence (104 occurrences), Execution (39), Discovery
(31), Not Malicious Yet (8), Impact (6), and Other (4). A command name is
therefore not sufficient to determine its tactic.

\textbf{Persistence example.} In the sequence
\texttt{echo root:<password> | chpasswd | bash}, the emitted string is consumed
by \texttt{chpasswd} to modify credentials. The command contributes to
maintaining access and is consequently labelled Persistence.

\textbf{Execution example.} In
\texttt{echo <encoded payload> | base64 --decode | bash}, the output is decoded
and passed directly to a shell. Here, \texttt{echo} participates in launching a
payload and is labelled Execution. These examples show that tactic assignment
depends on neighboring words and downstream use, motivating a contextual
sequence model rather than a command-to-label lookup table.

\subsection{Session Lengths} \label{subsec:task1_lengths}

\textbf{Length variability.} Training sessions range from 2 to 224 Bash words,
with mean \textbf{45.70} and median \textbf{17}. Test sessions have the same
range but are generally longer, with mean \textbf{57.38} and median
\textbf{23.5}. The large gap between means and medians indicates a
right-skewed distribution: many sessions are short, while a small number of
long traces extend the tail.

Figure~\ref{fig:task1_session_ecdf} visualizes the empirical cumulative
distribution over the labelled sessions. The long tail is relevant because
the number of Bash words does not directly equal the number of model tokens:
unusual shell strings can be fragmented into many subwords and exceed a
transformer's context limit even when the word count appears manageable.

\begin{figure}[tbp]
    \centering
    \includegraphics[width=0.92\linewidth]{\imagespath Task1/session_length_ecdf.png}
    \caption{ECDF of the number of Bash words per labelled SSH session in the
    training and test partitions, shown on linear and logarithmic scales.}
    \label{fig:task1_session_ecdf}
\end{figure}
